\begin{enabstract}

    3D light field display technology, as the core carrier of next-generation immersive visual experience, has broad application prospects in commercial advertising, medical diagnosis and treatment, film and television entertainment and other fields due to its advantage of glasses-free stereoscopic perception. As a key factor in constructing stereoscopic realism and spatial credibility, lighting editing quality directly determines the display effect and immersive experience of 3D light field portraits. However, existing 3D light field portrait lighting editing methods have technical pain points such as complex and expensive acquisition equipment, insufficient multi-view lighting consistency, and temporal flicker in dynamic scenes, making it difficult to meet the generation needs of high-quality static and dynamic light field content.
    
To address the above issues, this thesis conducts systematic research on 3D light field face lighting editing, constructing a complete technical system from static basic methods to dynamic optimization frameworks. The main research work and innovative achievements are as follows:

Firstly, a key information-driven 3D light field static face lighting editing method is proposed. Aiming at the problems of lighting inconsistency caused by independent multi-view editing and high cost of dense acquisition equipment in static scenes, a three-stage framework of "highlight removal-lighting editing-dense viewpoint generation" is designed: accurately locate and remove portrait highlight areas through a dual detection mechanism of brightness threshold and diffuse reflection component; adopt a dual-path lighting editing strategy that not only supports extracting Spherical Harmonic (SH) coefficients from HDR images to achieve environment light relighting, but also parses natural language instructions through BERT model to realize semantic regulation of lighting direction and intensity; generate 90 dense disparity maps within the horizontal ±30° range based on NeRF volume rendering technology to meet the needs of light field display. Experimental results show that the proposed method outperforms mainstream methods such as NVPR and LRPI in indicators including FID (33.55) and PSNR (32.06), with a single-frame inference time of only 25.92s. The vote rate of stereoscopic perception on 3D light field display is 89.19\%, and the vote rate of lighting consistency is 92.29\%.

Secondly, a time-space dual-dimension joint optimization framework for 3D light field face video dynamic lighting editing is proposed. Aiming at the temporal flicker and multi-view lighting mismatch problems when static methods are extended to video sequences, two core optimization modules are designed: the temporal consistency optimization module extracts stable semantic features through a shared encoder, combines optical flow alignment and lighting-geometry feature separation strategies, and applies dual loss constraints of "feature L1 + pixel residual L2" to effectively suppress inter-frame lighting flicker; the multi-view consistency optimization module achieves 3D geometric alignment based on camera internal and external parameters, converts multi-view lighting constraints into 3D point lighting observation problems through 3D surface sampling, and realizes physically consistent expression of cross-view lighting by combining a multi-factor aggregation mechanism with three weights: visibility, normal angle, and lighting confidence. Experimental verification shows that the framework reduces the temporal stability indicator WE to 0.0004 and the multi-view lighting error LE to 0.7922, with a comprehensive subjective experiment score of 4.73, which is significantly better than comparative methods such as PTI and SMFR.

Through theoretical innovation and experimental verification, this thesis effectively solves the core technical pain points in 3D light field face lighting editing, providing a new technical path for high-quality static and dynamic light field content generation. The relevant achievements can be applied in fields such as glasses-free 3D display, VR/AR, and film and television production, with important theoretical significance and engineering application value.

    \enkeywords{3D Light Field Display; Face Lighting Editing; Spherical Harmonic Function; NeRF; Temporal Consistency; Multi-view Consistency}
\end{enabstract}